% Chapter 7. Every number is a macro (generated/numbers.tex).
\chapter{Discussion and Limitations}
\label{ch:discussion}

\section{What the results say}
\label{sec:what}

\textbf{The silence belongs mostly to the recognizer.} On every utterance
length we measured, the final decode is the largest stage between the end of
speech and the first reply audio. The language model's first token is a small
stage. This answers RQ1. It also explains why the field's remedy, speculating
the reply, attacks the smaller target on this device. Our prediction that the
final's share would rise with length did not hold. The share is not monotone
in length and falls below half on two sets, so the claim stands as ``largest
stage everywhere'', not as a law of length.

\textbf{On shared silicon, speculation is never free.} RQ2's answer has three
parts:
\begin{itemize}
\item Co-residence alone slows the recognizer's final decode, through
  huge-page fallback (\TaxBallastA{} and \TaxBallastB{} in two sessions).
\item A reply draft still decoding at the endpoint costs \ThinkerOverlap.
\item A partial hypothesis still decoding at the endpoint costs \ListenerC\,ms
  of final decode per second of overlap.
\end{itemize}
The listener's penalty differs from the thinker's in one respect: it can be
removed at its source, by freezing the orphaned decode.

\textbf{The thinker's value is a timing problem.} RQ3's answer is a precise
negative. On spoken multi-step problems the question comes last, so a reply
draft becomes usable only near the end of the utterance. There, usability is
\FstarUsable{} at the last fraction where a draft still fits, within sampling
error of the break-even \BreakEven. The registered test at three-quarters of
the utterance closed the value side. Neither trigger evaluated could be shown
to predict the end. The negative is about this task's structure, not about
reply speculation in general.

\textbf{A simple rule is enough here.} RQ4's answer: with a feasibility gate
and a freeze, the listener cuts the median paired TTFA on long utterances by
\PthreeRed{} and abstains on short ones. The allocator never spends the window
on the thinker, and its decision log shows why. In simulation, a planner that
reasons about the endpoint's hazard improves on the gate's median by only
\PlannerDP. The rule is simple because, under the measured costs, little is
left to optimise at the median.

\section{Limitations}
\label{sec:limitations}

\textbf{One device, one model of each kind.} Every number comes from one
board class with CPU-bound recognition. A GPU recognizer would move the stage
shares, and possibly the balance between the consumers. The language model is
one small 4-bit model and the synthesizer one voice.

\textbf{Audio.} The benchmark audio is synthesized speech played in through
the capture path at real time; the development set is short assistant queries.
Live speech, from a single speaker, will be one descriptive pass that supports
no claim.

\textbf{The gate's prior is the set's own median.} The feasibility gate uses a
point prior derived from the set it runs on. A deployment would have to learn
the prior. On items longer than the prior the gate stops early, which is where
the planner's remaining gain sits. A prior that adapts to elapsed speech is an
untested hypothesis.

\textbf{The headline's interval.} The gated listener met its pre-registered
prediction, but the interval's lower end lies below the registered bar. The
ungated version failed the same prediction. The listener's overlap cost per
second lies below its registered band, and its scorer was written after the
data.

\textbf{Costs of the freeze.} Frozen decodes are resumed after the reply and
discarded, at a net mean of \StaleMean\,J each over \StaleN{} decodes. Some
drew far more power than their length explains, for reasons we have not
found. A turn that opens while the worker is frozen gets no hypotheses. The
freeze's result also depends on a quiet period after each reply.

\textbf{The occupancy tax.} The tax was measured in two attribution sessions,
with one run per condition made after its baseline, so time order is
confounded with condition. The sessions disagree on its size, and only the
second was instrumented for huge pages. The mechanism's explanation (buffer
churn against a fragmented pool) was offered after the fact. Only
per-process levers were tried against it.

\textbf{The thinker side.} Draft usability comes from an offline probe of 80
items per fraction, with no interval carried on the curve. The live
comparison of a continuing draft is descriptive: its interleaved design broke
at the first voice-activity split, and its scorer was written after the run.
The trigger evaluation is on the development set and is underpowered. Reading
the triggers on the multi-step set was registered not to happen once the
probe closed the value side.

\textbf{The planner is simulated.} It is calibrated from logs. Its
validation is weak: the final's cost per second was fitted largely from the
contrast the validation predicts, and the fill rate for unobserved commit
boundaries was chosen to match the measured committed fractions. Richer fills
raise the planner's median gain.

\textbf{Drift across dates.} Part of a cross-date difference in baseline TTFA
remains unattributed. Residency of the second recognizer was tested as a
cause and no cost was detected (\ResidencyTTFA). Every claim-bearing
comparison from pass 1n onward therefore has a same-day baseline; pass 2's
baseline ran on an earlier day, and its results are read with that caveat.

\section{Threats to validity}
\label{sec:threats}

\textbf{Internal.} Arms are paired on the item, within a day from pass 1n
onward, interleaved or blocked with a same-arm washout. Runs are checked against the file's ground
truth and voided by rules fixed before scoring. Runs on one day were voided
because the harness waited for a proxy instead of the event; those runs carry
no claim. One fence breach is on record, and its run is reported with and
without the affected repetition. Several scorers were written after the
data; each is marked in the ledger.

\textbf{Construct.} TTFA is measured from the true end of speech, known to the
audio source, so a late endpoint counts as latency. The onset-grid penalties
of the thinker are timed from speech onset, and are cited as differences
between arms on the same utterance. Usability is a word-prefix match of the
first synthesized chunk; a draft whose first chunk does not match counts as
unusable.

\textbf{External.} One device, one language model, one benchmark family.
The thinker's negative depends on questions coming last. The listener's gain
depends on long utterances and on a recognizer whose final decode dominates.
Both are stated as conditions, not as general results.

\textbf{Statistical.} Intervals are percentile bootstraps over items unless
stated. Some analyses resample turns or runs, and say so. Nulls are reported
with the effect sizes they exclude, and an underpowered null is reported as
``not detectable'', never as ``no effect''.
